% !TEX program = xelatex
% 编译: xelatex experiment_report.tex (需安装 TeX Live / MiKTeX 及 ctex)
\documentclass[12pt,a4paper]{ctexart}

\usepackage{amsmath,amssymb}
\usepackage{graphicx}
\usepackage{booktabs}
\usepackage{multirow}
\usepackage{geometry}
\usepackage{hyperref}
\usepackage{caption}
\usepackage{array}
\usepackage{float}

\geometry{left=2.5cm,right=2.5cm,top=2.5cm,bottom=2.5cm}
\hypersetup{colorlinks=true,linkcolor=blue,urlcolor=blue,citecolor=blue}

\title{\textbf{生物医学图像处理课程大作业}\\
\large Project 2：皮肤病变图像深度学习分类实验报告\\
\large 基于 ResNet18 与 MobileNetV2 的迁移学习}
\author{（请填写：组长姓名、学号、组员及分工）}
\date{\today}

\begin{document}
\maketitle

\begin{abstract}
本实验针对 Project 2 皮肤病变三分类任务（mel、nv、vasc），在预处理后的
\texttt{processed dataset} 上采用 ImageNet 预训练的 ResNet18 与 MobileNetV2 进行迁移学习。
训练集进一步按病灶原图 ID 分组划分验证集，在独立测试集（120 张）上评估，
并分别统计原图与增强图子集的表现。实验表明：MobileNetV2 在测试集整体准确率与
Macro-F1 上优于 ResNet18（83.33\% vs.\ 79.17\%，Macro-F1 0.840 vs.\ 0.797）；
在增强图子集（$n{=}91$）上 MobileNetV2 准确率为 79.1\%，高于 ResNet18（75.8\%）
及传统随机森林（71.4\%）。与 \texttt{final\_project} 手工特征分类器相比，
深度学习整体准确率提升约 10 个百分点（83.33\% vs.\ 73.33\%）。
\end{abstract}

\tableofcontents
\newpage

%==============================================================================
\section{实验背景与任务目标}
%==============================================================================

\subsection{课程任务说明}
Project 2 要求从皮肤病变图像中提取特征并完成多类分类，并通过对\textbf{增强图像}
的识别能力体现算法鲁棒性。数据集共 600 张图像：200 张原图与 400 张由数据增强
（\texttt{*\_aug1}、\texttt{*\_aug2}）得到的增强图，类别为：
\begin{itemize}
  \item \textbf{mel}：黑色素瘤（Melanoma）
  \item \textbf{nv}：痣（Nevus）
  \item \textbf{vasc}：血管病变（Vascular lesion）
\end{itemize}

\subsection{本实验目标}
\begin{enumerate}
  \item 使用 \textbf{ResNet18} 与 \textbf{MobileNetV2} 进行迁移学习分类；
  \item 完成训练集、验证集、测试集的划分与模型训练；
  \item 对比深度模型在\textbf{原图}与\textbf{增强图}上的分类表现，分析鲁棒性。
\end{enumerate}

\textit{说明：课程第 15 周定量评测要求使用传统图像处理或简单机器学习方法；
本报告所述深度学习结果为课堂展示与课程报告所用，与传统方法形成互补对比。}

%==============================================================================
\section{数据集与预处理}
%==============================================================================

\subsection{原始数据}
原始数据位于 \texttt{Data\_Proj2/}，包含 \texttt{image/}、\texttt{mask/} 及
\texttt{label.csv}（600 条记录）。

\subsection{预处理流程}
小组成员已完成统一预处理（见 \texttt{processed dataset/demo.ipynb}），主要步骤为：
\begin{enumerate}
  \item 尺寸归一化至 $256\times256$；
  \item 基于形态学 Black-hat 的轻度\textbf{去毛发}（Telea 修复）；
  \item LAB 色彩空间下对 L 通道做亮度归一化（范围约 30--220）。
\end{enumerate}

\subsection{划分方案}
预处理输出保存在 \texttt{processed dataset/}：
\begin{table}[H]
  \centering
  \caption{官方预处理划分（ stratify 分层抽样，random\_state=42）}
  \label{tab:data_split}
  \begin{tabular}{lrrrr}
    \toprule
    子集 & 样本数 & 原图 & 增强图 & 类别分布（nv/mel/vasc） \\
    \midrule
    训练（官方） & 480 & 171 & 309 & 240 / 168 / 72 \\
    测试（官方） & 120 & 29  & 91  & 60 / 42 / 18 \\
    合计         & 600 & 200 & 400 & 300 / 210 / 90 \\
    \bottomrule
  \end{tabular}
\end{table}

深度学习阶段在官方训练集（480 张）上，按\textbf{病灶原图编号}
（如 \texttt{42\_aug2} $\rightarrow$ \texttt{42}）进行 \textbf{GroupShuffleSplit}，
再划分出约 20\% 作为验证集，以避免同一病灶的原图与增强图同时出现在训练与验证中：
\begin{table}[H]
  \centering
  \caption{深度学习训练/验证划分（VAL\_RATIO=0.2，SEED=42）}
  \label{tab:dl_split}
  \begin{tabular}{lrrr}
    \toprule
    子集 & 样本数 & 原图 & 增强图 \\
    \midrule
    训练（模型拟合） & 387 & 138 & 249 \\
    验证（调参与早停） & 93  & 33  & 60  \\
    测试（最终评估）   & 120 & 29  & 91  \\
    \bottomrule
  \end{tabular}
\end{table}

%==============================================================================
\section{方法}
%==============================================================================

\subsection{迁移学习框架}
\begin{itemize}
  \item \textbf{骨干网络}：torchvision 提供的 ResNet18、MobileNetV2，权重为 ImageNet-1K 预训练；
  \item \textbf{分类头}：将最后一层全连接替换为 3 类输出；
  \item \textbf{输入}：$256\times256$ RGB，ImageNet 均值/方差归一化；
  \item \textbf{损失函数}：交叉熵（CrossEntropyLoss）；
  \item \textbf{优化器}：Adam（$lr=10^{-4}$，weight decay $10^{-4}$）；
  \item \textbf{学习率调度}：验证准确率 plateau 时乘以 0.5（patience=3）；
  \item \textbf{早停}：验证集准确率连续 8 个 epoch 无提升则停止。
\end{itemize}

\subsection{数据增强（仅训练阶段）}
训练时对预处理图像施加在线增强：随机水平翻转、$\pm15^\circ$ 旋转、
轻度 ColorJitter；验证与测试仅做 resize 与归一化。

\subsection{评估指标}
\begin{itemize}
  \item \textbf{Accuracy}（准确率）
  \item \textbf{Macro-F1}（宏平均 F1，兼顾类别不平衡）
  \item \textbf{混淆矩阵}及每类 Precision / Recall / F1
\end{itemize}

按 \texttt{image\_id} 是否包含 \texttt{\_aug} 将样本划分为\textbf{原图子集}
与\textbf{增强图子集}，分别统计上述指标。

%==============================================================================
\section{实验设置}
%==============================================================================

\begin{table}[H]
  \centering
  \caption{超参数与运行环境（\texttt{deep\_learning/config.py}）}
  \label{tab:hyper}
  \begin{tabular}{ll}
    \toprule
    项目 & 设置 \\
    \midrule
    批大小（batch size） & 16 \\
    最大训练轮数 & 30 \\
    验证集比例 & 0.2（按病灶分组） \\
    随机种子 & 42 \\
    实现框架 & PyTorch + torchvision \\
    代码入口 & \texttt{py -m deep\_learning.main} \\
    权重保存 & \texttt{outputs/checkpoints/\{backbone\}\_best.pt} \\
    结果输出 & \texttt{outputs/results/} \\
    \bottomrule
  \end{tabular}
\end{table}

%==============================================================================
\section{实验结果}
%==============================================================================

\subsection{总体性能对比}

表~\ref{tab:main_results} 汇总两模型在验证集与测试集上的准确率与 Macro-F1
（数据来源于 \texttt{outputs/results/model\_comparison.csv}）。

\begin{table}[H]
  \centering
  \caption{ResNet18 与 MobileNetV2 分类性能对比}
  \label{tab:main_results}
  \begin{tabular}{llrrrr}
    \toprule
    模型 & 数据集 & 子集 & $n$ & Accuracy & Macro-F1 \\
    \midrule
    \multirow{6}{*}{ResNet18}
      & val  & 全部   & 93  & 0.774 & 0.780 \\
      & val  & 原图   & 33  & 0.788 & 0.792 \\
      & val  & 增强图 & 60  & 0.767 & 0.774 \\
      & test & 全部   & 120 & 0.792 & 0.797 \\
      & test & 原图   & 29  & 0.897 & 0.899 \\
      & test & 增强图 & 91  & 0.758 & 0.760 \\
    \midrule
    \multirow{6}{*}{MobileNetV2}
      & val  & 全部   & 93  & 0.763 & 0.776 \\
      & val  & 原图   & 33  & 0.667 & 0.665 \\
      & val  & 增强图 & 60  & 0.817 & 0.833 \\
      & test & 全部   & 120 & \textbf{0.833} & \textbf{0.840} \\
      & test & 原图   & 29  & 0.966 & 0.971 \\
      & test & 增强图 & 91  & \textbf{0.791} & \textbf{0.793} \\
    \bottomrule
  \end{tabular}
\end{table}

\subsection{测试集：原图与增强图性能差距}

\begin{table}[H]
  \centering
  \caption{测试集上原图 vs.\ 增强图的准确率差距（$\Delta$ = 原图 Acc $-$ 增强图 Acc）}
  \label{tab:gap}
  \begin{tabular}{lrrr}
    \toprule
    模型 & 原图 Acc & 增强图 Acc & $\Delta$ \\
    \midrule
    ResNet18     & 89.7\% & 75.8\% & +13.9 pp \\
    MobileNetV2  & 96.6\% & 79.1\% & +17.5 pp \\
    \bottomrule
  \end{tabular}
\end{table}

ResNet18 在测试原图上表现突出（89.7\%），增强图回落至 75.8\%（差距约 14 pp）。
MobileNetV2 测试原图准确率达 96.6\%（$n{=}29$，样本量小，可能受划分泄漏影响偏乐观），
增强图为 79.1\%，为两深度模型中\textbf{增强子集最高}；整体测试 Macro-F1 为 0.840，
优于 ResNet18 的 0.797。

\subsection{测试集分类报告（MobileNetV2，最佳模型）}

\begin{table}[H]
  \centering
  \caption{MobileNetV2 测试集整体 per-class 指标}
  \label{tab:mb_test_report}
  \begin{tabular}{lrrrr}
    \toprule
    类别 & Precision & Recall & F1-score & Support \\
    \midrule
    mel  & 0.85 & 0.79 & 0.81 & 42 \\
    nv   & 0.83 & 0.83 & 0.83 & 60 \\
    vasc & 0.81 & 0.94 & 0.87 & 18 \\
    \midrule
    accuracy & \multicolumn{3}{r}{0.83} & 120 \\
    macro avg & 0.83 & 0.85 & 0.84 & 120 \\
    \bottomrule
  \end{tabular}
\end{table}

\begin{table}[H]
  \centering
  \caption{ResNet18 测试集整体 per-class 指标}
  \label{tab:rn_test_report}
  \begin{tabular}{lrrrr}
    \toprule
    类别 & Precision & Recall & F1-score & Support \\
    \midrule
    mel  & 0.81 & 0.69 & 0.74 & 42 \\
    nv   & 0.78 & 0.83 & 0.81 & 60 \\
    vasc & 0.80 & 0.89 & 0.84 & 18 \\
    \midrule
    accuracy & \multicolumn{3}{r}{0.79} & 120 \\
    macro avg & 0.80 & 0.80 & 0.80 & 120 \\
    \bottomrule
  \end{tabular}
\end{table}

ResNet18 对 \textbf{mel} 的召回率偏低（0.69），是拉低增强子集表现的主要因素之一；
MobileNetV2 在三类上更为均衡，\textbf{vasc} 召回率达 0.94。

\subsection{混淆矩阵}

混淆矩阵行表示真实类别，列表示预测类别，顺序均为 mel / nv / vasc。

\begin{table}[H]
  \centering
  \caption{MobileNetV2 测试集混淆矩阵（$n=120$）}
  \label{tab:cm_mb_test}
  \begin{tabular}{l|ccc}
    \toprule
    真实 $\backslash$ 预测 & mel & nv & vasc \\
    \midrule
    mel  & 33 & 9  & 0 \\
    nv   & 6  & 50 & 4 \\
    vasc & 0  & 1  & 17 \\
    \bottomrule
  \end{tabular}
\end{table}

\begin{table}[H]
  \centering
  \caption{ResNet18 测试集混淆矩阵（$n=120$）}
  \label{tab:cm_rn_test}
  \begin{tabular}{l|ccc}
    \toprule
    真实 $\backslash$ 预测 & mel & nv & vasc \\
    \midrule
    mel  & 29 & 12 & 1 \\
    nv   & 7  & 50 & 3 \\
    vasc & 0  & 2  & 16 \\
    \bottomrule
  \end{tabular}
\end{table}

\begin{table}[H]
  \centering
  \caption{MobileNetV2 测试集——增强图子集混淆矩阵（$n=91$）}
  \label{tab:cm_mb_aug}
  \begin{tabular}{l|ccc}
    \toprule
    真实 $\backslash$ 预测 & mel & nv & vasc \\
    \midrule
    mel  & 22 & 8  & 0 \\
    nv   & 6  & 39 & 4 \\
    vasc & 0  & 1  & 11 \\
    \bottomrule
  \end{tabular}
\end{table}

\begin{table}[H]
  \centering
  \caption{ResNet18 测试集——增强图子集混淆矩阵（$n=91$）}
  \label{tab:cm_rn_aug}
  \begin{tabular}{l|ccc}
    \toprule
    真实 $\backslash$ 预测 & mel & nv & vasc \\
    \midrule
    mel  & 19 & 10 & 1 \\
    nv   & 7  & 39 & 3 \\
    vasc & 0  & 1  & 11 \\
    \bottomrule
  \end{tabular}
\end{table}

主要误分类模式：mel 与 nv 之间的相互混淆（两模型均有 7--15 例量级）；
ResNet18 将部分 mel 误判为 vasc（1 例），增强子集上 mel$\rightarrow$nv 误判更为突出。

\subsection{训练曲线与错例可视化}

图~\ref{fig:curve_rn}--\ref{fig:curve_mb} 为两模型训练/验证损失及准确率、Macro-F1 曲线；
图~\ref{fig:errors_mb}--\ref{fig:errors_rn} 为测试集部分错例（真标签 $\rightarrow$ 预测标签）。
由 \texttt{py -m deep\_learning.main --plot} 根据
\texttt{*\_history.json} 与预测 CSV 生成，保存在
\texttt{deep\_learning/outputs/figures/}。

\begin{figure}[H]
  \centering
  \includegraphics[width=0.92\linewidth]{deep_learning/outputs/figures/resnet18_training_curves.png}
  \caption{ResNet18 训练曲线（损失与准确率 / Macro-F1）}
  \label{fig:curve_rn}
\end{figure}

\begin{figure}[H]
  \centering
  \includegraphics[width=0.92\linewidth]{deep_learning/outputs/figures/mobilenet_v2_training_curves.png}
  \caption{MobileNetV2 训练曲线}
  \label{fig:curve_mb}
\end{figure}

MobileNetV2 错例中 mel 误判为 nv 仍占多数；ResNet18 在增强样本上错例更多。

\begin{figure}[H]
  \centering
  \includegraphics[width=0.92\linewidth]{deep_learning/outputs/figures/mobilenet_v2_test_misclassified.png}
  \caption{MobileNetV2 测试集错例示意（优先展示 mel 相关误判）}
  \label{fig:errors_mb}
\end{figure}

\begin{figure}[H]
  \centering
  \includegraphics[width=0.92\linewidth]{deep_learning/outputs/figures/resnet18_test_misclassified.png}
  \caption{ResNet18 测试集错例示意}
  \label{fig:errors_rn}
\end{figure}

\subsection{训练策略改进}

后续训练默认启用：\textbf{类别加权交叉熵}（缓解 nv/mel/vasc 不平衡）、
以验证集 \textbf{Macro-F1} 选取最佳 checkpoint（而非单纯准确率）。
复现命令见附录；推理输出 \texttt{output\_\{backbone\}.csv} 可由
\texttt{--predict} 生成。

%==============================================================================
\section{结果分析与讨论}
%==============================================================================

\subsection{模型对比}
在独立测试集上，\textbf{MobileNetV2 整体优于 ResNet18}（准确率 83.33\% vs.\ 79.17\%，
Macro-F1 0.840 vs.\ 0.797）。尽管 ResNet18 参数量更大，在本任务有限样本（600 张）
及已预处理数据上，轻量级 MobileNetV2 配合迁移学习取得了更好泛化；在增强测试子集上
亦以 79.1\% 高于 ResNet18 的 75.8\%。

\subsection{原图与增强图鲁棒性}
Project 2 强调对增强图像的识别鲁棒性。实验表明：
\begin{itemize}
  \item \textbf{ResNet18}：测试集原图准确率（89.7\%）高于增强图（75.8\%），差距约 14 pp，
  对离线增强（旋转、翻转）泛化不足；
  \item \textbf{MobileNetV2}：增强图准确率 79.1\%，为深度模型中最高；原图准确率达 96.6\%
  （$n{=}29$），与原图--增强图差距约 17.5 pp，部分源于原图子集样本少及划分泄漏带来的乐观估计，
  应以增强子集指标衡量鲁棒性。
\end{itemize}

验证集上 MobileNetV2 呈现“增强优于原图”的趋势，与测试集原图子集高准确率并存，
与验证/测试原图样本量均较小（$n{=}33$ / $29$）有关，最终以增强子集与整体测试集为准。

\subsection{类别不平衡}
数据集中 nv 样本最多，vasc 最少。MobileNetV2 对 vasc 召回率达 0.94、F1 为 0.87；
ResNet18 对 mel 召回为 0.69，临床意义下 melanoma 漏检风险相对更高，可结合
代价敏感学习或重采样进一步优化。

\subsection{局限性}
\begin{enumerate}
  \item 官方 train/test 按单张 \texttt{image\_id} 随机划分，不同增强版本可能分属训练与测试，
  存在病灶级信息泄漏的可能；本实验仅在 train$\rightarrow$val 划分时按病灶分组缓解；
  \item 测试集增强图占比高（91/120），评测结果更侧重增强分布；
  \item 与传统方法的对比见第~\ref{sec:trad_vs_dl} 节；深度学习未复现成员4的全部
  图像扰动实验，鲁棒性结论主要基于原图/增强图子集划分。
\end{enumerate}

%==============================================================================
\section{传统方法 vs 深度学习对比分析}
\label{sec:trad_vs_dl}
%==============================================================================

本节基于 \texttt{final\_project/} 中小组已完成的传统流水线，与
\texttt{deep\_learning/} 实验结果进行\textbf{同数据集、同官方划分}
（训练 480 / 测试 120，\texttt{random\_state=42}）的并列对比。
传统方法数据来自 \texttt{results/model\_comparison.csv}、
\texttt{results\_member4/augmentation\_consistency.csv} 与
\texttt{results\_member4/group\_split\_summary.csv}；
深度学习方法数据来自 \texttt{deep\_learning/outputs/results/model\_comparison.csv}。

\subsection{传统方法流水线概述}

传统方案由小组成员分模块完成（详见 \texttt{final\_project/README.md}）：
\begin{enumerate}
  \item \textbf{预处理}（成员1）：与深度学习共用同一套
  \texttt{processed dataset}（$256\times256$、去毛发、亮度归一化）；
  \item \textbf{特征工程}（成员2）：从 RGB 图像提取 127 维手工特征——
  颜色矩（9 维）、RGB 直方图（96 维）、GLCM 纹理（12 维）、LBP 纹理（10 维），
  经 \texttt{StandardScaler} 归一化（可选 PCA 降至 36 维，最终评测使用 127 维）；
  \item \textbf{分类器}（成员3）：对 SVM（RBF）、随机森林（RF）、KNN 进行
  \texttt{GridSearchCV}（5 折分层交叉验证），在测试集上报告 Accuracy 与 Macro-F1；
  \item \textbf{鲁棒性扩展}（成员4）：分析增强图一致性、图像扰动及严格按病灶
  编号分组划分下的性能波动。
\end{enumerate}

\subsection{总体性能对比}

表~\ref{tab:trad_vs_dl_overall} 汇总各方法在官方测试集（$n=120$）上的整体指标。
传统方法中 \textbf{随机森林（RF）} 准确率最高（73.33\%），与 SVM（72.50\%）Macro-F1
相当（均为 0.736）；KNN 明显落后（66.67\%）。
深度学习中 \textbf{MobileNetV2} 全面领先：准确率 83.33\%、Macro-F1 0.840，
较最佳传统方法 RF 分别提升约 \textbf{10.0 pp} 与 \textbf{0.104}；
ResNet18 亦优于全部传统模型（79.17\% / 0.797）。

\begin{table}[H]
  \centering
  \caption{传统机器学习 vs.\ 深度学习：测试集整体性能（$n=120$）}
  \label{tab:trad_vs_dl_overall}
  \begin{tabular}{llrrl}
    \toprule
    范式 & 模型 & Accuracy & Macro-F1 & 备注 \\
    \midrule
    \multirow{3}{*}{传统（手工特征）}
      & Random Forest & 0.733 & 0.736 & 最佳传统准确率 \\
      & SVM (RBF)     & 0.725 & 0.736 & 少数类 vasc 召回略优 \\
      & KNN           & 0.667 & 0.672 & $k{=}1$，曼哈顿距离 \\
    \midrule
    \multirow{2}{*}{深度学习（迁移学习）}
      & MobileNetV2   & \textbf{0.833} & \textbf{0.840} & 本实验最佳 \\
      & ResNet18      & 0.792 & 0.797 & 参数量更大 \\
    \bottomrule
  \end{tabular}
\end{table}

\subsection{原图与增强图子集对比}

Project 2 强调对增强图像的识别能力。测试集中含原图 29 张、增强图 91 张
（\texttt{*\_aug1}、\texttt{*\_aug2}）。表~\ref{tab:trad_vs_dl_aug} 对比各方法在
两类子集上的准确率及\textbf{原图--增强图差距} $\Delta$。

\begin{table}[H]
  \centering
  \caption{原图 vs.\ 增强图测试准确率及差距 $\Delta$（pp = 百分点）}
  \label{tab:trad_vs_dl_aug}
  \begin{tabular}{lrrrr}
    \toprule
    模型 & 原图 Acc ($n{=}29$) & 增强图 Acc ($n{=}91$) & $\Delta$ & 说明 \\
    \midrule
    SVM (传统)           & 89.7\% & 67.0\% & +22.7 & 增强图表现最差 \\
    Random Forest (传统) & 79.3\% & 71.4\% & +7.9  & 传统中最均衡 \\
    KNN (传统)           & 75.9\% & 63.7\% & +12.2 & 对增强图敏感 \\
  \midrule
    ResNet18             & 89.7\% & 75.8\% & +13.9 & 原图高、增强回落明显 \\
    MobileNetV2          & 96.6\% & 79.1\% & +17.5 & \textbf{增强子集最佳} \\
    \bottomrule
  \end{tabular}
\end{table}

\textit{注：传统方法增强子集准确率为 \texttt{augmentation\_consistency.csv} 中
\texttt{\_aug1}（45 张）与 \texttt{\_aug2}（46 张）子集准确率的样本数加权平均。}

主要发现如下：
\begin{itemize}
  \item \textbf{绝对性能}：MobileNetV2 在增强图子集（79.1\%）上较最佳传统 RF（71.4\%）
  高出约 8 pp，亦高于 ResNet18（75.8\%）；
  \item \textbf{泛化均衡性}：传统 RF 的原图--增强差距（7.9 pp）最小；深度模型原图
  准确率普遍偏高，MobileNetV2 原图达 96.6\%（$n{=}29$），提示需结合严格病灶分组
  评估，不宜单独依据原图子集判断鲁棒性；
  \item \textbf{共同难点}：mel 与 nv 互混仍是主要误分类来源（见错例图）。
\end{itemize}

\subsection{各类别 per-class 表现}

表~\ref{tab:trad_vs_dl_class} 对比最佳传统模型 RF 与最佳深度模型 MobileNetV2
的测试集各类 F1。两者共同难点为 \textbf{mel $\leftrightarrow$ nv} 混淆；
深度学习在三类上均优于传统方法，尤其对 mel（0.81 vs.\ 0.71）与 vasc（0.87 vs.\ 0.76）
提升明显。深度 MobileNetV2 将 mel 召回提升至 0.79，有利于降低黑色素瘤漏检风险。

\begin{table}[H]
  \centering
  \caption{测试集各类 F1-score：Random Forest（传统）vs.\ MobileNetV2（深度）}
  \label{tab:trad_vs_dl_class}
  \begin{tabular}{lrrr}
    \toprule
    类别 & RF (传统) & MobileNetV2 & $\Delta$ F1 \\
    \midrule
    mel  & 0.707 & 0.810 & +0.103 \\
    nv   & 0.744 & 0.830 & +0.086 \\
    vasc & 0.757 & 0.870 & +0.113 \\
    \midrule
    macro avg & 0.736 & 0.840 & +0.104 \\
    \bottomrule
  \end{tabular}
\end{table}

\subsection{数据划分与评测可信度}

成员4 指出：当前官方划分按单张 \texttt{image\_id} 随机分层，去除
\texttt{\_aug1}/\texttt{\_aug2} 后缀后，训练集与测试集存在
\textbf{98 个共同病灶编号}（共 200 个原图编号），即同一病灶的原图/增强图可能
分属训练与测试，造成信息泄漏，使测试指标偏乐观。

在\textbf{严格按病灶编号分组}（10 次随机，训练/测试病灶零重叠）下，
传统方法 Macro-F1 均值显著下降（表~\ref{tab:strict_group}）。
深度学习仅在 train$\rightarrow$val 阶段采用 GroupShuffleSplit 缓解泄漏，
测试集仍使用官方划分，故与传统“官方划分”结果可直接对比，但双方均需在答辩中
说明严格分组下的预期性能折损。

\begin{table}[H]
  \centering
  \caption{严格病灶分组（10 次均值$\pm$标准差）下的 Macro-F1}
  \label{tab:strict_group}
  \begin{tabular}{lcc}
    \toprule
    模型 & Macro-F1 (mean$\pm$std) & 相对官方划分 RF（0.736） \\
    \midrule
    Random Forest (传统) & $0.578 \pm 0.035$ & $-0.158$ \\
    SVM (传统)           & $0.515 \pm 0.050$ & $-0.221$ \\
    KNN (传统)           & $0.490 \pm 0.067$ & $-0.246$ \\
    \bottomrule
  \end{tabular}
\end{table}

此外，成员4 对同一测试集施加亮度、对比度、高斯噪声、模糊、旋转裁剪等扰动：
传统 SVM 在亮度 $\times 0.8$ 时准确率降至 31.7\%，强噪声（$\sigma{=}16$）时降至 25\%；
RF 相对稳定但仍出现 5--15 pp 量级的回落。深度学习流水线未系统复现上述扰动实验，
但 MobileNetV2 在\textbf{课程定义的离线增强图}上保持 79.1\% 准确率，仍高于传统 RF（71.4\%）
与 ResNet18（75.8\%）。

\subsection{方法特性与适用场景}

\begin{table}[H]
  \centering
  \caption{传统方法 vs.\ 深度学习：综合对比}
  \label{tab:trad_vs_dl_summary}
  \small
  \begin{tabular}{p{2.2cm}p{5.2cm}p{5.2cm}}
    \toprule
    维度 & 传统（手工特征 + SVM/RF/KNN） & 深度学习（ResNet18 / MobileNetV2） \\
    \midrule
    特征表示 &
    可解释的颜色/纹理统计量，127 维，依赖领域设计 &
    ImageNet 预训练卷积自动表征，端到端微调 \\
    训练成本 &
    特征提取一次 + 网格搜索约 1--2 分钟 &
    GPU 训练数十分钟，需调参（学习率、早停等） \\
    官方划分性能 &
    最佳 73.33\%（RF），满足课程第 15 周定量要求 &
    最佳 83.33\%（MobileNetV2），适合课堂展示 \\
    增强图鲁棒性 &
    RF 增强子集 71.4\%，SVM 原图--增强差距大 &
    MobileNetV2 增强 79.1\%，整体最优 \\
    严格分组泛化 &
    Macro-F1 约 0.49--0.58，波动较大 &
    未在测试集做严格分组；val 已按病灶分组 \\
    临床相关 &
    mel 召回偏低（RF 0.69） &
    mel 召回 0.79，vasc 召回 0.94 \\
    \bottomrule
  \end{tabular}
\end{table}

\textbf{结论性认识}：在相同预处理与划分条件下，深度学习以更高上限完成三分类任务，
MobileNetV2 为综合最优；传统 RF/SVM 实现简单、可解释、训练快，符合课程对传统
方法的考核，但在有限样本皮肤镜任务上难以匹敌迁移学习。课程设计中的“传统 +
深度”双轨互补：前者展示特征工程与经典分类器基线，后者展示表示学习与鲁棒性优势。
后续可在严格病灶分组测试集上统一评测深度模型，并探索“手工特征 + CNN 嵌入”融合。

%==============================================================================
\section{结论}
%==============================================================================

本实验完成了基于 ResNet18 与 MobileNetV2 的皮肤病变三分类迁移学习流水线，
在预处理数据集上实现了训练、验证、测试全流程，并按原图/增强图分别评估。
\textbf{MobileNetV2} 在测试集取得最佳整体性能（Accuracy 83.33\%，Macro-F1 0.840），
增强图子集准确率 79.1\%，为所对比方法中最高。
\textbf{ResNet18} 测试准确率 79.17\%，对增强图泛化弱于 MobileNetV2，mel 召回为 0.69。

与 \texttt{final\_project} 传统流水线对比：最佳传统模型 RF 测试准确率 73.33\%、
Macro-F1 0.736，增强子集约 71.4\%；MobileNetV2 整体提升约 10 pp、增强子集提升约 8 pp。
严格病灶分组下传统 Macro-F1 降至约 0.49--0.58，提示需在无泄漏划分下进一步验证深度模型。
课程要求的双路径——传统方法满足定量基线，深度学习提供更高精度与增强鲁棒性。

已实现类别加权损失、Macro-F1 选模、训练曲线与错例可视化及 \texttt{output.csv} 推理脚本；
后续可在严格分组测试集上复评深度模型，并尝试手工特征与 CNN 嵌入融合。

%==============================================================================
\section*{参考文献}
\addcontentsline{toc}{section}{参考文献}
\begin{enumerate}
  \item He K, Zhang X, Ren S, Sun J. Deep Residual Learning for Image Recognition. CVPR, 2016.
  \item Sandler M, Howard A, Zhu M, et al. MobileNetV2: Inverted Residuals and Linear Bottlenecks. CVPR, 2018.
  \item 课程 Project 2 说明文档：\texttt{Project idea.pdf}
  \item 实验代码与结果：\texttt{deep\_learning/}，评估结果 \texttt{deep\_learning/outputs/results/}
  \item 传统方法流水线：\texttt{final\_project/}，对比数据 \texttt{results/}、\texttt{results\_member4/}
\end{enumerate}

%==============================================================================
\appendix
\section{复现实验命令}
%==============================================================================

\begin{verbatim}
cd 大作业目录
py -m pip install -r deep_learning/requirements.txt
py -m deep_learning.main --stats
py -m deep_learning.main --train --eval --plot --backbone all
py -m deep_learning.main --predict --backbone mobilenet_v2
py -m deep_learning.main --predict --submitdemo --backbone mobilenet_v2
\end{verbatim}

输出文件包括：\texttt{model\_comparison.csv}、
\texttt{*\_history.json}、\texttt{outputs/figures/} 下图、
\texttt{output\_mobilenet\_v2.csv}，以及各划分的预测结果 CSV。
关闭类别权重：加 \texttt{--no-class-weights}。

\section{小组成员分工（请填写）}
\begin{table}[H]
  \centering
  \begin{tabular}{lll}
    \toprule
    姓名 & 学号 & 主要贡献 \\
    \midrule
    & & 数据预处理 \\
    & & 深度学习模型训练与评估 \\
    & & 实验报告撰写 \\
    \bottomrule
  \end{tabular}
\end{table}

\end{document}
